\section{问题一模型建立及求解}

% v8.0.0：一级标题由 Template Authority 固定；以下二级标题只是复杂问题的默认占位，
% 可依据当前数学对象改成更专业的题目专属标题。MODEL / SOLVE / RESULT / VALIDATE
% 是功能槽，不是必须逐字出现的固定小节名。简单解析题可合并或省略非必要槽位。

\subsection{模型建立}
% 从本题对象和机制开始，详细推导关键关系。核心公式应说明来源、推导关键理由和后续用途。
设……，得到
\begin{equation}
    F(x,\theta)=0,
    \label{eq:q1-core-relation}
\end{equation}
式~\eqref{eq:q1-core-relation} 控制……，并用于下一步构造……。

% 命题只在确实承担降维、等价变换、单调性、可行性等建模作用时使用。
% \begin{hskproposition}{变换模型的等价性}
% \label{prop:q1-equivalence}
% 在条件……及参数范围……内，原模型与变换模型具有相同的可行解映射和最优目标值。
% \begin{hskproof}
% 由定义构造……，并验证约束与目标在映射前后保持一致，故命题得证。
% \end{hskproof}
% \end{hskproposition}

% 核心模型收束是 rendering mode，不默认建立“核心模型汇总”独立二级标题。
% 对复杂优化模型，在模型建立末尾先单独展示 objective，再单独展示 constraints；
% 目标函数不得放入约束大括号。简单解析问题可只在正文中收束最终关系，不强制展示本段。
\begin{equation}
    \min_{\mathbf{x}} f(\mathbf{x})
    \label{eq:q1-objective}
\end{equation}

\begin{equation}
\text{s.t.}\quad
\left\{
\begin{aligned}
    g_i(\mathbf{x}) &\le 0, \quad i=1,\ldots,m, \\
    h_j(\mathbf{x}) &= 0, \quad j=1,\ldots,n, \\
    \mathbf{x} &\in \Omega .
\end{aligned}
\right.
\label{eq:q1-constraints}
\end{equation}
% 非优化题按真实结构改成状态方程、观测关系、概率映射、边界/初值或最终解析关系；
% 不为了模板对称保留无意义的 objective / s.t. 示例。

\subsection{模型求解}
% 先说明最终模型被化成什么可计算问题以及剩余困难，再说明 solver 为什么适配；
% 写真实初始化/参数、精度或停止条件和输出映射，不要只写“利用 Python/MATLAB 求解”。

\subsection{求解结果}
% 主结果、方案和核心数值放在这里。正文核心图/表必须有邻近的显式交叉引用和本问解释。
% 最后一个结果段自然回答本问，不默认再建“问题一结论”。

% 当存在实质敏感性、鲁棒性、多初值/多 seed、边界扰动、替代模型/算法、残差/区间或外样本证据时，
% 可启用下列功能槽；若只是简单解析或直接计算，不为结构对称强制增加。
% \subsection{结果的分析与验证}
% 先说明需要检查的风险，再说明扰动/比较方式、评价指标、结果变化以及主结论可支持到什么范围。

% 按实际题目复制本文件为 07_question2.tex、08_question3.tex 等，并只写各问新增结构。
% 共享基础模型只完整定义一次，后问不从头复制。
